\documentclass{article}
\usepackage{/globals/math}
\usepackage{/globals/de}
\begin{document}
\section{Dreisatz} 
Mit dem Dreisatz können die Werte von Unbekannten in Sachkontexten mit zwischen zwei Paare an zwei Variablen, welche alle in bestimmer Relation zueinander stehen beantwortet werden.

Es werden sich zwei Paare gedacht, ein Paar aus $a$ und $b$ und ein Paar aus $x$ und $y$. Ein Paar besteht aus zwei sachkontexlichen Werten, die in Relation zueinander Stehen (\zb \textquote{Kilogramm Mehl zu Anzahl gebackener Brötchen} oder \textquote{Anzahl MitarbeiterInnen zu benötigte Zeit}). Ist nun der Wert einer der Variablen unbekannt, so können die Relationen der anderen Variablen genutzt werden um den Wert der Unbekannten zu finden.
Hier wird unterschieden zwischen zwei Arten an Relation
\begin{description}
 \item[Proportionalität] Wächst der Wert der einen Variable, so wächst auch der Wert der anderen Variable (\textquote{je mehr, desto mehr}). Hier gilt
\[
 \frac{x}{y} = \frac{a}{b}
\]
 \item[Antiproportionalität] Wächst der Wert der einen Variable, so schrinkt der Wert der anderen Variable (\textquote{je mehr, desto weniger}). Hier gilt
\[
 x*y=a*b
\]
\end{description}
Es ist hilfreich sich die Einheiten der Variable anzusehen und zu überlegen, ob diese Sinn ergeben. Diese Verhalten sich den allgemeinen algebraischen Regeln; sie können bei einer Division wegfallen, so dass das Ergebnis auch die Einheit hat, welche erwartet wird.

Ansonsten kann frei nach Intuition gerechnet werden.
\end{document}